\section{Selmer data and the descent}\label{sec:selmer}

The denominator bound of \cref{thm:denominator-bound} asks that the twists by $\vartheta^{-1}_\ell$ of the main summands have short Selmer groups. This section proves that boundedly
many auxiliary primes, \emph{tied} so that tame inertia acts like Frobenius times a \emph{slope}, make the tied Selmer groups of finitely many lattices short
(\cref{thm:descent}). The next section turns the ties into a twist.

\subsection{Selmer data}\label{ss:selmer-data}
Let $E$ be a CM field in $\Omega$, $\ell$ a prime and $O$ the integers of a finite extension of $\Q_\ell$. A \emph{Selmer datum} for a finite set $S_0$ of places and an element $g_0$ is a lattice $T=O^n$ with a Galois action $\rho$,
unramified outside a set $S\subseteq S_0$, a module $Z_w$ of local cocycles at each $w\mid\ell$, and an element $\gamma_*$ acting by a scalar $\zeta_*$ with $\zeta_*^2-1$ a unit, such that $\det(\rho(g)-1)\ne0$ for $g=g_0$ and $g=g_0^{-c}=(g_0^c)^{-1}$.
The \emph{local relations} are the linear relations $\sum_ja_jz(g_j)=0$, $g_j$ in a decomposition group above $\ell$, that hold on the $Z_w$. A class in $H^1(E,T/a)$ is identified with its cocycle vanishing at $\gamma_*$. The \emph{global term}
is the supremum over $a\ne0$ of the lengths of the modules of equivariant maps from $T/a$ to its dual twisted by complex conjugation, the \emph{places term} the sum of the local terms at the places of $S$ prime to $\ell$, and
$\beta_T$ is the sum of these and of the valuations of $\det(\rho(g)-1)$ for $g=g_0,g_0^{-c}$.
\lean{New.Selmer.conjDualHom, New.Selmer.congruenceTerm, New.Selmer.localTerm, New.Selmer.placeTerms, New.Selmer.firstPrimeTerms, New.Selmer.SelmerDatum, New.Selmer.SelmerDatum.globalTerm, New.Selmer.SelmerDatum.placesTerm, New.Selmer.SelmerDatum.firstPrimeTerm, New.Selmer.SelmerDatum.localTerms, New.Selmer.SelmerDatum.β, New.Selmer.SelmerDatum.β_ne_top, New.Selmer.SelmerDatum.isUnit_ζ_sub_one, New.Selmer.SelmerDatum.localRelations, New.Selmer.SelmerDatum.mem_obeying_localRelations_iff, New.Selmer.SelmerDatum.coboundary_mem_obeying_localRelations, New.Selmer.SelmerDatum.length_localCocycles_le_finrank_mul, New.Selmer.SelmerDatum.classRelations, New.Selmer.SelmerDatum.tiedSelmer, New.Selmer.SelmerDatum.apply_γ_eq_zero, New.Selmer.evalConjPair, New.Selmer.evalConjPair_apply, New.Selmer.mem_selmerZ1_of_mem_selmerModule, New.Selmer.IsNewAuxiliaryPrime, New.Selmer.IsTiedPrimes.character_tameGenerator_eq_ite, New.Selmer.IsTiedPrimes.insert_of_isNewAuxiliaryPrime, New.Selmer.smul_mem_span_of_forall_smul_mem, New.Selmer.eq_top_of_stable_of_isUnit, New.Selmer.SelmerDatum.ρ_γ, New.Selmer.SelmerDatum.ρ_conjGal_γ_inv, New.Selmer.SelmerDatum.cyclotomicChar_γ, New.Selmer.SelmerDatum.sahGroup, New.Selmer.SelmerDatum.commutator_γ_mem, New.Selmer.SelmerDatum.commutator_conjGal_γ_mem, New.Selmer.SelmerDatum.sahGroup_le, New.Selmer.SelmerDatum.conjGal_mem_sahGroup, New.Selmer.SelmerDatum.map_eq_one_of_mem_sahGroup, New.Selmer.SelmerDatum.ρ_eq_one_of_mem_sahGroup, New.Selmer.SelmerDatum.sahGroup_le_normalClosure, New.Selmer.SelmerDatum.bijective_γ_smul_sub, New.Selmer.SelmerDatum.cocycle_conj_of_fixed, New.Selmer.SelmerDatum.cocycle_commutator_γ, New.Selmer.SelmerDatum.cocycle_eq_zero_of_sahGroup, New.Selmer.SelmerDatum.exists_mem_sahGroup_isUnit_apply, New.Selmer.SelmerDatum.span_image_sahGroup_eq_top, New.Selmer.length_torsionBy_pi_quot_le, New.Selmer.span_range_smul_single_eq_top, New.Selmer.length_torsionBy_characterModule_le, New.Selmer.smul_apply_eq_zero_of_mem_conjDualHom, New.Selmer.congruenceTerm_le}

Write $W=T/a=(O/a)^n$, $G=\Gal(\Omega/E)$, $g^c$ for the conjugate of $g$ by complex conjugation and $g^{-c}=(g^c)^{-1}$; $W^*$ is the character module $\Hom(W,\Q/\Z)$ with $g$ acting by $y\mapsto\mathrm{cyc}(g)\,y\circ g^{-1}$, and $W^{*c}$ is $W^*$ with $g$ acting as $g^c$. A Selmer datum also asks that $T/\varpi$ be absolutely
irreducible in Burnside's form (the matrices $\rho(g)$ span all matrices), that $\gamma_*$ be inertial at a place of $S$ prime to $\ell$, and that $\sum_{w\mid\ell}(\rank Z_w-n)\le fn$, $f$ the number of infinite places. The group $H$ is generated, as an abstract group, by the commutators $[g,\gamma_*]$ and
$[g,\gamma_*^c]$, $g\in G$.

\begin{lemma}[Classes on the commutator group]\label{lem:selmer-basics}
Let $T$ be a Selmer datum and $a\ne0$.
\begin{enumerate}[label=\textup{(\alph*)}]
\item Coboundaries obey the local relations, and at each $w\mid\ell$ the restrictions of the cocycles obeying them have length at most $\rank Z_w\cdot\length(O/a)$.
\item $H$ is normal and stable under $g\mapsto g^c$; homomorphisms to commutative groups and $\rho$ are trivial on $H$; and $H$ lies in the normal subgroup generated by the inertia groups above any set containing the places of $S$ prime to $\ell$ and their conjugates.
\item (Sah's identity.) A cocycle $z$ with $z(\gamma_*)=0$ has $z([g,\gamma_*])=(1-\zeta_*)z(g)$, and $z(ghg^{-1})=gz(h)$ when $h$ acts as the identity. A cocycle that vanishes at $\gamma_*$ and on $H$ vanishes.
\item If a submodule of $W$ is stable under a set of matrices spanning all matrices and contains a vector with a unit coordinate, it is $W$.
\item If $g$ acts by a scalar $\alpha$ and $g^{-c}$ by $\alpha'$, then $\Delta=\alpha-\mathrm{cyc}(g^c)\alpha'$ kills $\Hom_G(W,W^{*c})$, and the $\Delta$-torsion of $W^*$ is no longer than that of $W$, which has length at most $n\,v(\Delta)$.
\end{enumerate}
\end{lemma}

\begin{proof}
(a) The coboundary of $b\in W$ lifts to the coboundary of a lift of $b$ in $O^n$, which lies in each $Z_w$, so it obeys every relation that holds on $Z_w$. For the bound, fix a coordinate $i$ and let $\phi(f)$ be the linear form $z\mapsto(\sum_gf(g)z(g))_i$ on $Z_w$, for finitely supported families $f$ of matrices. A function $c$ obeying the relations of $Z_w$ gives the map $f\mapsto(\sum_gf(g)c(g))_i$ to $O/a$, which kills $\ker\phi$ (row $i$ of a family in $\ker\phi$, the other
rows set to zero, is a relation on $Z_w$) and determines $c$ (take for $f$ the families with one nonzero entry). Hence the module embeds into the $O$-linear maps from $\mathrm{im}\,\phi$ to $O/a$, and $\mathrm{im}\,\phi$ lies in the dual of $Z_w$, free of rank at most $\rank Z_w$.

(b) $g'[g,\delta]g'^{-1}=[g'g,\delta][g',\delta]^{-1}$ shows normality, and $[g,\gamma_*]^c=[g^c,\gamma_*^c]$ stability under conjugation. A homomorphism to a commutative group kills commutators. As $\gamma_*$ and $\gamma_*^c$ act by scalars, $\rho$ kills the commutators. Finally $\gamma_*$ and $\gamma_*^c$ lie in inertia groups above
such a set, hence in its normal closure, and so does every commutator with them.

(c) By the cocycle rule, $z(\gamma_*^{-1})=-\gamma_*^{-1}z(\gamma_*)=0$ and $z(g^{-1})=-g^{-1}z(g)$, so $z(\gamma_*g^{-1}\gamma_*^{-1})=\zeta_*z(g^{-1})$ and $z(g\gamma_*g^{-1}\gamma_*^{-1})=z(g)-\zeta_*z(g)$. If $h$ acts as the identity, $z(ghg^{-1})=z(g)+gz(h)-ghg^{-1}z(g)=gz(h)$. If $z$ vanishes on $H$, then $(1-\zeta_*)z(g)=0$, and $1-\zeta_*$ is a unit as
$\zeta_*^2-1$ is.

(d) If $u$ has a unit coordinate $u_j$, the matrix whose only nonzero column is the $j$th, a lift of $m/u_j$, takes $u$ to $m$; such matrices are combinations of the given ones.

(e) An equivariant $\eta$ satisfies $\eta(gu)=g\cdot\eta(u)$; as $g$ acts on $W$ by $\alpha$ and on $W^{*c}$ by $\mathrm{cyc}(g^c)\alpha'$ (the action of $g^c$ on $W^*$ is by $\mathrm{cyc}(g^c)$ times the transpose inverse of $\rho(g^c)$, which is $\alpha'$), $\Delta\eta=0$. A character killed by $\Delta$ is a character of $W/\Delta W$, which has as many elements as
$W[\Delta]$; and $W[\Delta]=\prod_j(O/a)[\Delta]$ has length at most $n\,v(\Delta)$.
\end{proof}

\begin{lemma}[Sah]\label{lem:sah}
Let $H$ be the group generated by the commutators $[g,\gamma_*]$ and $[g,\gamma_*^c]$. If a cocycle vanishing at $\gamma_*$ has a value with a unit coordinate, its values on $H$ span $T/a$. If $g$ acts by a scalar $\alpha$, $g^{-c}$ by $\alpha'$ and the image of
$\rho$ spans all matrices, the global term is at most $n\,v(\alpha-\mathrm{cyc}(g^c)\alpha')$.
\end{lemma}

\begin{proof}
For the first claim, let $z(g)$ have a unit coordinate. By \cref{lem:selmer-basics}(c), $z([g,\gamma_*])=(1-\zeta_*)z(g)$ has a unit coordinate too, and $[g,\gamma_*]\in H$. The span of $z(H)$ is stable under $G$, since $gz(h)=z(ghg^{-1})$ for $h\in H$, which acts as the identity, and $H$ is normal (\cref{lem:selmer-basics}(b)); by
\cref{lem:selmer-basics}(d) it is $T/a$. For the second claim, by \cref{lem:selmer-basics}(e) $\Delta=\alpha-\mathrm{cyc}(g^c)\alpha'$ kills every equivariant map $\eta\colon W\to W^{*c}$. Evaluation at a basis vector $e$ is injective on these maps: a map vanishing at $e$ vanishes at its translates, which span $W$ because the image of
$\rho$ spans all matrices. Hence $\Hom_G(W,W^{*c})$ embeds into the $\Delta$-torsion of $W^*$, of length at most $n\,v(\Delta)$; take the supremum over $a$.
\end{proof}

\subsection{Auxiliary primes and ties}\label{ss:aux-primes}
\emph{Auxiliary prime data} at a prime $P\ne\bar P$ of norm $\equiv1$ modulo $q$, outside a set $T$, consist of a place of $\Omega$ above $P$, a Frobenius element $\varphi$, the chosen tame generator $\gamma$, and a character to $\Z/q$ that comes
from an idele class character of level $P$ trivial at $T$ and takes the value $1$ at $\gamma$. A set $J$ of \emph{tied primes} consists of a first prime $P_0$, whose Frobenius element lies in a given set $W$, and later primes, all
with auxiliary prime data for $q=\ell^N$ and, for each datum $i$, a \emph{slope} $s_v^{(i)}\in O$, with $s_{P_0}=0$, such that at the primes of nonzero slope for $i$ the Frobenius element and its conjugate inverse satisfy a condition
$V_i$ (typically: they act as the identity on $T_i/\ell^M$). The \emph{ties} of datum $i$ are $z(\gamma_v)=s_vz(\varphi_v)$ and $z(\gamma_v^{-c})=s_vz(\varphi_v^{-c})$, $v\in J$. The \emph{tied Selmer group} $\mathrm{Sel}_a(J)$ modulo $a$ consists of the cocycles with values
in $T/a$, vanishing at $\gamma_*$, unramified outside $S\cup J\cup\bar J$, obeying the local relations and the ties. A prime $P$ with data is \emph{new auxiliary} for datum $i$ if it may be added to $J$ with any slope for $i$.
\lean{New.Selmer.AuxPrimeData, New.IsAuxPrimeDataAt, New.IsAuxiliary, New.Selmer.exists_discreteLog, New.Selmer.valuation_pow_sub_one_lt_one_of_isFrob, New.Selmer.exists_ideleClassChar_of_local, New.Selmer.exists_anticyclotomicChar_of_frob_fixes_roots, New.Selmer.mem_ker_inf_stabilizer_of_mem_normalClosure, New.Selmer.exists_isAuxPrimeDataAt, New.Selmer.reduceModPow, New.Selmer.IsTiedPrimes, New.Selmer.IsTiedPrimes.apply_idOrConjInv_of_slope_ne_zero, New.Selmer.tieRelations, New.Selmer.idOrConjInv_mem_inertiaAbove}
\lean{New.Selmer.affineAction_apply, New.Selmer.affineAction_pow_apply, New.Selmer.affineAction_pow_eq_one, New.Selmer.commute_affineAction, New.Selmer.sub_one_smul_eq_of_commute, New.Selmer.inv_sub_one_smul_eq, New.Selmer.sub_one_smul_eq_of_tame, New.Selmer.lineDirection, New.Selmer.selmerOnLine, New.Selmer.shiftedDet, New.Selmer.eval₂_shiftedDet, New.Selmer.coeff_shiftedDet_card, New.Selmer.tiePolynomial, New.Selmer.selmerAtPoint_zero, New.Selmer.selmerOnLine_zero, New.Selmer.map_smul_eq_of_reduceModPow_eq, New.Selmer.dvd_eval_universalChar_sub_one, New.Selmer.exists_monoidHom_smul_eq_pointTwist_smul, New.Selmer.sub_one_smul_eq_of_isAuxPrimeDataAt, New.Selmer.IsTiedPrimes.frob_sub_one_smul_eq, New.Selmer.IsTiedPrimes.frob_sub_one_smul_tie_eq_zero_of_selmerOnLine, New.Selmer.IsTiedPrimes.relLength_tied_le_sum_ker, New.Selmer.IsTiedPrimes.relLength_tied_le, New.Selmer.IsTiedPrimes.relLength_tieRelations_le_sum_addVal_det}

A prime of $J$ is in particular \emph{auxiliary} for $S_0$: of prime norm, with neither it nor its conjugate in $S_0$, and not conjugate to another prime of $J$. A tied set also asks that for $v\ne P_0$ the character of $P_0$ take the value $1$ at $\varphi_v$ and the characters of the other primes the value $0$, and that no prime of $S_0$ or of $J$ divide the
norm of $P_0$. For characters $h_j$ to $\Z/q$ and a point $x$ the \emph{twist} by $\psi_x(g)=\prod_j(1+x_j)^{h_j(g)}$ is the action $g\mapsto\psi_x(g)\rho(g)$ (\cref{ss:universal-char}); $\mathrm{Sel}(x)$ is its Selmer module.

\begin{lemma}[Characters at auxiliary primes]\label{lem:aux-characters}
Let $E$ be a CM field and $q=\ell^N$.
\begin{enumerate}[label=\textup{(\alph*)}]
\item (Discrete logarithm.) If $\Nm P=1+qk$ is prime and the residue of a local unit $u$ generates the residue units, some character of $E_P^\times$ to $\Z/q$ sends $u$ to $1$ and kills every $t$ with $t\equiv1$ or $t^k\equiv1$ modulo $P$.
\item If a Frobenius element $\sigma$ at a place above $P\notin T$, $\Nm P=1+nk$, fixes an $n$th root of a $T$-unit $u$, then $u^k\equiv1$ modulo $P$.
\item If the $T$-ideles and the principal ideles generate the finite ideles, a continuous character of $E_P^\times$ killing the $T$-units extends to a character of the idele classes that is, on the $T$-ideles, the character of the component at $P$.
\item (Anticyclotomic characters.) Under the hypothesis of (c), if $P\notin T$, $P\ne\bar P$, $\Nm P=1+qk$ is prime and a Frobenius $\sigma$ at a place $A$ over $P$ fixes $q$th roots of generators of the $T$-units, there is a continuous character $\chi\colon G\to\Z/q$ with value $1$ at the tame generator at $A$
and $\chi(\mathrm{Art}(x))=\theta(x/\bar x)$ for a character $\theta$ of the idele classes trivial on the $T$-ideles congruent to $1$ at $P$.
\item If $q$ is a unit off $T$ and the places of $R$ lie off $T$, the normal closure of the inertia groups above $R$ fixes the $q$th roots of unity and $q$th roots of $T$-units.
\end{enumerate}
\end{lemma}

\begin{proof}
(a) Divide out the power of a uniformiser and take the residue; the character of the residue units, cyclic of order $qk$, that sends the generator to $1$ kills the $k$th powers' preimages of $1$: if the residue of $t$ is $g^e$ with $g^{ek}=1$, the order $qk$ of $g$ divides $ek$, so $q\mid e$.

(b) With $y^n=u$, $y=\sigma y\equiv y^{\Nm P}=y\,u^k$ modulo the place, and $y$ is a unit there.

(c) The $T$-ideles map onto the idele classes with kernel the $T$-units, so the character of the component at $P$, which kills the $T$-units, descends; it is continuous as it agrees on the open group of $T$-ideles with a continuous character.

(d) By (b) the discrete logarithm of (a), for a generator $u_0$ of the residue units, kills the $T$-units, and by (c) it extends to an idele class character $\theta$; its kernel is open, containing all powers with exponent its finite index. The existence theorem of class field theory (\cref{thm:B-classfield}(c)) turns $x\mapsto\theta(x/\bar x)$ into a
Galois character $\chi$. The tame generator at $A$ is the Artin symbol of $u_0$ at $P$, whose conjugate has component $1$ at $P$ as $P\ne\bar P$, so $\chi$ takes the value $\theta(u_0)=1$ there.

(e) The inertia groups at places off $T$, and their conjugates, fix the roots of unity of order $q$ (prime to these places) and the $q$th roots of the $T$-units, which are units there; the fixer of these finitely many roots is a normal open subgroup.
\end{proof}

\lean{New.Selmer.unionWithConj, New.Selmer.pointChar, New.Selmer.pointChar_apply}

\begin{lemma}[Ties and the line]\label{lem:untied}
Let $J$ be tied primes for $S_0$, $i$ a datum, and let $\rho$ modulo $\ell^M$ be trivial on inertia outside $S_0$ and at the elements with $V_i$.
\begin{enumerate}[label=\textup{(\alph*)}]
\item (Affine actions.) For a cocycle $c$, $g\mapsto(u\mapsto c(g)+gu)$ is a homomorphism to the affine group. If $g$ acts by a scalar $\alpha$ with $\sum_{i<k}\alpha^i=0$, the $k$th power of its affine map is the identity. If the affine maps of $\varphi$ and $\gamma$ commute, then $(\varphi-1)c(\gamma)=(\gamma-1)c(\varphi)$, and the relation passes
from $\varphi$ to $\varphi^{-1}$ if $\gamma$ acts by a scalar.
\item (Twists.) If $y$ divides all coordinates of $x$, then $y\mid\psi_x(g)-1$. If the $1+x_j$ are $q$th roots of unity modulo $\ell^M$, the twisted matrices act on $(O'/\ell^M)^n$ through a representation $R$, and $R(g)$ is the scalar $\psi_x(g)$ wherever $\rho(g)\equiv1$.
\item (The relation at a point of the line.) Let $x_{P_0}=y$, $x_v=s_vy$ for $v\ne P_0$, and let $\sum_{j<q}\alpha^j\equiv0$ modulo $\ell^M$ whenever $\alpha\equiv1$ modulo $y$. For $z\in\mathrm{Sel}(x)$ and $w$ a prime $v\in J$ or its conjugate, the matrix $(1+y)\rho(\varphi_w)-1$ kills $z(\gamma_w)-s_vz(\varphi_w)$; for $v=P_0$ this needs
$y=0$.
\item The quotient of $\mathrm{Sel}(x)$ by its tied part on $I\subseteq J$ (with $P_0\in I$ only if $y=0$) has length at most $\sum_w\length\ker((1+y)\rho(\varphi_w)-1)$, hence at most $v(P^1(y))$ for the \emph{tie polynomial} $P^1(Y)=\prod_w\det((1+Y)\rho(\varphi_w)-1)$ over $w\in(J\setminus P_0)\cup\overline{(J\setminus P_0)}$.
\item (Release of the first prime.) If $M\le N$ and $\varphi_{P_0}$, $\varphi_{P_0}^{-c}$ act modulo $\ell^M$ as $g_0$, $g_0^{-c}$, adding the relations $z(\gamma_{P_0})=0=z(\gamma_{P_0}^{-c})$ to the Selmer module of $\rho$ modulo $\ell^M$ lowers its length by at most $\sum_{g=g_0,g_0^{-c}}v(\det(\rho(g)-1))$.
\end{enumerate}
\end{lemma}

\begin{proof}
(a) The homomorphism property is the cocycle rule. If $g$ acts by $\alpha$, the $k$th iterate is $u\mapsto\alpha^ku+\sum_{i<k}\alpha^ic(g)$, and $\sum_{i<k}\alpha^i=0$ gives $\alpha^k=1$ as $(\alpha-1)\sum\alpha^i=\alpha^k-1$. Commuting affine maps give $c(\gamma\varphi)=c(\varphi\gamma)$, which by the cocycle rule is
$c(\gamma)+\gamma c(\varphi)=c(\varphi)+\varphi c(\gamma)$. For $\varphi^{-1}$, multiply by $-\varphi^{-1}$ and use $c(\varphi^{-1})=-\varphi^{-1}c(\varphi)$ and that $\gamma$ commutes with $\varphi^{-1}$ on $T/a$.

(b) Each factor $(1+x_j)^{h}-1$ is divisible by $x_j$, hence by $y$, and a product of elements $\equiv1$ modulo $y$ is $\equiv1$. Modulo $\ell^M$, $\psi_x$ is a character $\chi$ (the $1+x_j$ being $q$th roots of unity and $h_j$ having values in $\Z/q$), and $R=\chi\rho$.

(c) Modulo $\ell^M$ the twist is the representation $R$ of (b). Inertia above $S_0$ acts through $R$ by scalars $\equiv1$ modulo $y$, and as the identity near $1$; so each inertial element acts by $\alpha$ with $\sum_{j<q}\alpha^j\equiv0$, and its affine map has order dividing $q$ by (a). As $q$ is prime to the place $v$ (its norm is
$\equiv1$ modulo $q$) and the cocycle is locally constant, the affine maps of inertia at $v$ form a finite group of exponent dividing $q$ on which wild inertia acts as the identity; conjugation by $\varphi_v$ is the $\Nm v$th power on tame inertia, and $\Nm v\equiv1$ modulo $q$, so the affine maps of $\varphi_v$ and $\gamma_v$
commute, and (a) gives $(R(\varphi_w)-1)z(\gamma_w)=(R(\gamma_w)-1)z(\varphi_w)$ at $v$ and, by the same argument at the conjugate place, at $\bar v$. Moreover $\lambda_u(\gamma_v)=[u=v]$ and, for $v\ne P_0$, $\lambda_u(\varphi_v)=[u=P_0]$, so $\psi_x(\gamma_v)=1+s_vy$ and $\psi_x(\varphi_v)=1+y$: the relation reads
$((1+y)\rho(\varphi_w)-1)z(\gamma_w)=s_vyz(\varphi_w)$. If $s_v\ne0$, $V_i$ gives $\rho(\varphi_w)\equiv1$ modulo $\ell^M$, so $s_v((1+y)\rho(\varphi_w)-1)z(\varphi_w)=s_vyz(\varphi_w)$ as well, and the difference is the claim; if $s_v=0$ the claim is the relation itself. At $P_0$ with $y=0$ the point $x$ is $0$ and both
values of $\psi$ are $1$.

(d) The map $z\mapsto(z(\gamma_w)-s_vz(\varphi_w))_w$ has the tied part as kernel and, by (c), its image in $\bigoplus_w\ker((1+y)\rho(\varphi_w)-1)$. The kernel of a matrix $A$ on $(O'/\ell^M)^n$ has length at most $v(\det A)$ (Smith normal form), and the determinant of $(1+y)\rho(\varphi_w)-1$ is the value at $y$ of
$\det((1+Y)\rho(\varphi_w)-1)$, a polynomial with leading coefficient $\det\rho(\varphi_w)$.

(e) At $y=0$ the point $x$ vanishes and $\mathrm{Sel}(0)$ is the Selmer module of $\rho$; by (c) the two tie maps take values in the kernels of $\rho(\varphi_{P_0})-1$ and $\rho(\varphi_{P_0}^{-c})-1$, which act modulo $\ell^M$ as $g_0-1$ and $g_0^{-c}-1$, and (d) concludes.
\end{proof}

\begin{proposition}[Auxiliary primes]\label{prop:aux-primes}
\begin{enumerate}[label=\textup{(\alph*)}]
\item (Chebotarev.) For finite sets $S$, $R$ and a neighbourhood $W$ of an element of the normal closure of the inertia groups at the places of $R$ outside $S\cup\bar S$ and not above $\ell$, there is a prime $P\notin R\cup\bar R$ of prime norm with
auxiliary prime data whose Frobenius element lies in $W$ and is killed by the character.
\item The character of $w\in J$ takes the value $1$ at the tame generator of $v$ if $w=v$ and $0$ otherwise.
\item (The relation from tame inertia.) If $\varphi$ is a Frobenius element of norm $\equiv1$ modulo $q$ and $\gamma$ inertial, each inertial element acting by a scalar $\alpha$ with $\sum_{i<q}\alpha^i=0$, then every cocycle has
$(\varphi-1)c(\gamma)=(\gamma-1)c(\varphi)$.
\item Slope $0$ means unramified: if all slopes of datum $i$ vanish, the tied Selmer group lies in the Selmer group without auxiliary primes; a cocycle that vanishes at $\gamma$, $\gamma^{-c}$ of a new auxiliary prime is unramified there.
\end{enumerate}
\end{proposition}

\begin{proof}
(a) Take a finite set $T$ of places containing $S$, $\bar S$ and the places above $\ell$, but no other place of $R$, whose ideles together with the principal ideles generate the finite ideles. The $T$-units are finitely generated; let $U$ be the open subgroup fixing the $q$th roots of unity and $q$th roots of a set of generators.
By \cref{lem:aux-characters}(e), $U$ contains the given element $z$, as the relevant inertia groups sit at places off $T$. Chebotarev's theorem (\cref{thm:B-classfield}(d)) gives $P\ne\bar P$ off $T\cup R\cup\bar R$ of prime norm and a Frobenius $\sigma$ at a place $A$ over $P$ such that every Frobenius at $A$ lies in $U$ and has $W$ as a
neighbourhood. As $\sigma$ fixes the roots of unity of order $q$, $\Nm P\equiv1$ modulo $q$; \cref{lem:aux-characters}(d) gives the character $\chi$, with value $1$ at the tame generator; multiplying $\sigma$ by a power of the tame generator gives a Frobenius $\tau\in W$ with $\chi(\tau)=0$. An idele that is a unit off a set
$T'\subseteq T$ closed under conjugation and $\equiv1$ at $P$ is killed by $\theta$, which gives the auxiliary prime data for every $T'\subseteq S\cup\{v\mid\ell\}$.

(b) The character of $w$ takes the value $1$ at its own tame generator by construction. For $v\ne w$ in $J$, $v$ is neither $w$ nor $\bar w$, and the character of $w$ comes from an idele class character of level $w$, unramified outside $w$, $\bar w$; so it kills the inertia at $v$, in particular the tame generator.

(c) The affine maps of inertia at the place form a finite group: the cocycle is locally constant and $\rho$ is trivial on inertia near $1$. Each has $q$th power $1$ by \cref{lem:untied}(a), as it acts by a scalar $\alpha$ with $\sum_{i<q}\alpha^i=0$; $q$ is prime to the place, so wild inertia acts as the identity and the group is a quotient of tame inertia, on
which $\varphi$ acts by the $\Nm$th power, $\Nm\equiv1$ modulo $q$. Hence the affine maps of $\varphi$ and $\gamma$ commute, and \cref{lem:untied}(a) gives the relation, also for $\varphi^{-1}$.

(d) At a prime $v\in J$ outside $S_0$, inertia acts as the identity on $T/a$, so a cocycle restricts on inertia to a homomorphism; it factors through tame inertia, as wild inertia is pro-$p$ with $p\ne\ell$ and the values are $\ell$-power torsion, and the tame quotient is topologically generated by $\gamma_v$. If all slopes of datum $i$ vanish, the
ties read $z(\gamma_v)=0=z(\gamma_v^{-c})$, so $z$ vanishes on inertia at the chosen places above $v$ and $\bar v$, and at the other places above them by $z(g\tau g^{-1})=gz(\tau)$ (\cref{lem:selmer-basics}(c)): $z$ is unramified at $J\cup\bar J$. The same argument applies at a new auxiliary prime.
\end{proof}

\subsection{The count}\label{ss:selmer-count}
\lean{New.Selmer.SelmerDatum.smul_eq_self_of_reduceModPow_eq_one, New.Selmer.SelmerDatum.isLocallyConstant_smul, New.Selmer.SelmerDatum.reductionModule, New.Selmer.SelmerDatum.reductionKer, New.Selmer.SelmerDatum.mem_reductionKer, New.Selmer.SelmerDatum.smul_eq_self_of_mem_reductionKer, New.Selmer.SelmerDatum.mem_reductionKer_of_mem_sahGroup, New.Selmer.SelmerDatum.mem_reductionKer_of_mem_inertiaAbove, New.Selmer.SelmerDatum.mem_obeying_tieRelations_iff, New.Selmer.SelmerDatum.coboundary_mem_obeying_tieRelations, New.Selmer.SelmerDatum.coboundary_mem_selmerModule, New.Selmer.SelmerDatum.apply_first_tameGenerator_eq_zero_of_mem_tiedSelmer, New.Selmer.SelmerDatum.mem_selmerModule_erase, New.Selmer.SelmerDatum.tiedSetting, New.Selmer.SelmerDatum.mem_cocycles_tiedSetting, New.Selmer.SelmerDatum.tiedSetting_avoids, New.Selmer.SelmerDatum.dualSelmer, New.Selmer.SelmerDatum.exists_length_localCocycles_tied_le, New.Selmer.SelmerDatum.length_add_length_tiedSelmer_le, New.Selmer.SelmerDatum.exists_length_cocycles_le, New.Selmer.SelmerDatum.sum_length_localCocycles_placesDividing_le, New.Selmer.SelmerDatum.sum_length_localCocycles_placeOver_le, New.Selmer.SelmerDatum.exists_sum_length_localCocycles_tied_le, New.Selmer.SelmerDatum.exists_length_tiedSelmer_le}
\lean{New.Selmer.isTameGeneratorFor_conjGal_inv, New.Selmer.characterModule_add_apply, New.Selmer.characterModule_zero_apply, New.Selmer.characterModule_eq_zero_of_span, New.Selmer.SelmerDatum.reductionKerInfConj, New.Selmer.SelmerDatum.mem_reductionKerInfConj, New.Selmer.SelmerDatum.isOpen_reductionKerInfConj, New.Selmer.SelmerDatum.sahGroup_le_reductionKerInfConj, New.Selmer.SelmerDatum.mem_reductionKerInfConj_of_mem_inertiaGroup, New.Selmer.SelmerDatum.γ₀_mem_reductionKerInfConj, New.Selmer.SelmerDatum.γ₀_mul_mem_normalClosure_inertiaAbove, New.Selmer.SelmerDatum.exists_isNewAuxiliaryPrime, New.Selmer.SelmerDatum.exists_apply_pair_eq_nsmul, New.Selmer.SelmerDatum.mirrorForm, New.Selmer.SelmerDatum.frobInTameImage, New.Selmer.SelmerDatum.apply_conj_of_mem_dualSelmer, New.Selmer.SelmerDatum.apply_first_tameGenerator_eq_zero_of_mem_dualSelmer, New.Selmer.SelmerDatum.eq_zero_of_mem_dualSelmer_of_forall_sahGroup, New.Selmer.SelmerDatum.dualSelmer_frob_add_conj_frob_eq_zero, New.Selmer.SelmerDatum.mirrorForm_self_eq_zero, New.Selmer.SelmerDatum.mirrorForm_add_swap_eq_zero, New.Selmer.SelmerDatum.mirrorForm_conj_γ_left, New.Selmer.SelmerDatum.mirrorForm_conj_γ_right, New.Selmer.SelmerDatum.mirrorForm_eq_zero, New.Selmer.SelmerDatum.exists_linearMap_apply_eq_dualSelmer_conj, New.Selmer.SelmerDatum.mem_conjDualHom_of_apply_eq, New.Selmer.SelmerDatum.length_dualSelmer_le}

Let $U$ be the kernel of the action on $T/\ell^M$ and on the roots of unity of order $\ell^N$, and $U\cap U^c$ the set of the $g$ with $g,g^c\in U$; for $a\mid\ell^M$, $W$ is a discrete Galois module, its orbit maps factoring through $T/\ell^M$. The \emph{tied setting} is the global duality setting of \cref{thm:B-greenberg-wiles} for
$W$, the primes $S'=S\cup I\cup\bar I$ with $I=J\setminus\{P_0\}$, chosen places above them, the kernel $U$, and at each $v\in S'$ the local cocycles that are restrictions of cocycles obeying the local relations and the ties. Its dual $\mathrm{Sel}^*_a(J)$, for a finite set $R$, is the module of the equivariant homomorphisms $U\to W^*$,
extended by $0$, unramified outside $R\setminus\{P_0,\bar P_0\}$ and obeying reciprocity at the split primes outside it.

\begin{lemma}[The tied setting]\label{lem:tied-setting}
Let $J$ be tied primes with $V_i$ the triviality on $T/\ell^M$, $M<N$ and $a\mid\ell^M$.
\begin{enumerate}[label=\textup{(\alph*)}]
\item $U$ contains $H$ and the inertia groups at the primes outside $S$; coboundaries obey the ties and lie in the Selmer module; tied classes vanish at $\gamma_{P_0}$, $\gamma_{P_0}^{-c}$ and are unramified at $P_0$, $\bar P_0$.
\item At a prime $v\in J$ outside $S$, or its conjugate, the local cocycles obeying the local relations and the ties are determined by their value at $\varphi_v$ (or $\varphi_v^{-c}$), so form a module no longer than $W$.
\item The local modules above $\ell$ have total length at most $(\#\{w\mid\ell\}+f)\length W$, those at the places of $S$ at most $(\#S+f)\length W$ plus the places term, and those at all of $S'$ at most $(\#S'+f)\length W$ plus the places term.
\item $\length W+\length\mathrm{Sel}_a(J)\le\length Z(\emptyset)$, where $Z(\emptyset)$ is the module of cocycles unramified outside $S'$ with the local conditions of the setting.
\item (Greenberg--Wiles.) For a suitable finite set $R$, $\length Z(\emptyset)+(\#S'+f)\length W\le\length\mathrm{Sel}^*_a(J)+\length W+\sum_{v\in S'}\length L_v$, $L_v$ the local modules.
\end{enumerate}
\end{lemma}

\begin{proof}
(a) $H$ acts as the identity through $\rho$ and fixes the roots of unity, as the cyclotomic character is a homomorphism to a commutative group (\cref{lem:selmer-basics}(b)); inertia at primes off $S$ acts as the identity on $T$, and fixes the roots of unity of order $\ell^N$ when the prime is prime to $\ell$. The tame generators and, where the slope is not zero, the
Frobenius elements fix $W$, so the coboundary of $b$ has $(\gamma_v-1)b=0=s_v(\varphi_v-1)b$; with \cref{lem:selmer-basics}(a) it lies in the Selmer module. As $s_{P_0}=0$, the ties at $P_0$ read $z(\gamma_{P_0})=0=z(\gamma_{P_0}^{-c})$, and \cref{prop:aux-primes}(d) makes $z$ unramified there; so $P_0$ can be removed from the
allowed ramification.

(b) A tied cocycle that vanishes at the Frobenius element vanishes at the tame generator by the tie, hence on inertia (\cref{prop:aux-primes}(d)) and on the decomposition group, which Frobenius and inertia generate topologically.

(c) By \cref{lem:selmer-basics}(a) the module at $w\mid\ell$ has length at most $\rank Z_w\cdot\length(O/a)$, and $\sum_w\rank Z_w\le(\#\{w\mid\ell\}+f)n$. At a place of $S$ prime to $\ell$ all classes are allowed, and the local term bounds the excess over $\length W$. On $S$ the ties only shrink the modules; above $I\cup\bar I$ use (b).

(d) Map $W\times\mathrm{Sel}_a(J)\to Z(\emptyset)$ by $(b,z)\mapsto z+\partial b$. It lands in $Z(\emptyset)$ by (a), and it is injective: $z(\gamma_*)=0$ and $(\partial b)(\gamma_*)=(\zeta_*-1)b$, with $\zeta_*-1$ a unit.

(e) This is \cref{thm:B-greenberg-wiles}(a) for the tied setting. Its hypothesis holds: $\gamma_*$ fixes the roots of unity (it is inertial at a place prime to $\ell$), is central modulo $U$ (it acts on $T/\ell^M$ by a scalar), and fixes no nonzero vector, as $\zeta_*-1$ is a unit. The norm of $P_0$ can be avoided by the auxiliary primes of the count:
it is prime to $\ell^N$, no prime of $S'$ divides it, and the inertia groups above its prime divisors lie in $U$. Writing both sides as orders of modules and passing to lengths gives the inequality.
\end{proof}

\begin{lemma}[Reciprocity at conjugate primes]\label{lem:mirror}
In the situation of \cref{lem:tied-setting}, let $\gamma_0$ be the tame generator at $P_0$.
\begin{enumerate}[label=\textup{(\alph*)}]
\item $U\cap U^c$ is open and contains $H$, $\gamma_0$ and the inertia groups at the primes outside $S_0$; for $h\in H$, $\gamma_0h$ lies in the normal closure of the inertia groups at the places of a set $R\supseteq\{P_0\}\cup S_0\cup\bar S_0$ prime to $\ell$ and not among the later primes of $J$ or their conjugates.
\item For $h\in H$, a finite set $R$ and a neighbourhood $X$ of $\gamma_0h$, there is a new auxiliary prime $P\notin R\cup\bar R$ for datum $i$ whose Frobenius element lies in $X$.
\item The conjugate inverse $\delta^{-c}$ of a normalised tame generator $\delta$ at $B$ is a normalised tame generator at $\bar B$, and the values of a cocycle at an inertial $\tau$, $\tau^{-c}$ are the same multiple of its values at $\delta$, $\delta^{-c}$.
\item A member $y$ of the dual is equivariant on $U$, vanishes at $\gamma_0$ and $\gamma_0^c$, and vanishes if it vanishes on $H$.
\item (Reciprocity.) If $P$ is a new auxiliary prime with $P,\bar P\notin R$ and $\varphi,\varphi^c\in U$, every $z$ in the tied Selmer group with ramification allowed also at $P$, $\bar P$ has $y(\varphi)(tz(\gamma))+y(\varphi^c)(tz(\gamma^{-c}))=0$ for $t\in O$.
\end{enumerate}
\end{lemma}

\begin{proof}
(a) $U$ is open and conjugation is continuous; $H\subseteq U$ by \cref{lem:tied-setting}(a) and $H^c=H$; inertia at a prime outside $S_0$, and its conjugate, act as the identity. The element $\gamma_0$ is inertial at $P_0$, and $h$ is a product of commutators with $\gamma_*$, $\gamma_*^c$, which lie in inertia groups at places of $S_0\cup\bar S_0$ prime to
$\ell$ (\cref{lem:selmer-basics}(b)).

(b) Apply \cref{prop:aux-primes}(a) with the set $R$ enlarged by the primes dividing the norm of $P_0$, the later primes of $J$ and their conjugates, and with a neighbourhood of $\gamma_0h$ inside $X\cap U\cap U^c$ on which the characters of the primes of $J$ are constant. Characters are trivial on $H$, so their values there are those at
$\gamma_0$: $1$ for $P_0$ and $0$ for the others, as required of a new prime; and $V_i$ holds at $\varphi$ and $\varphi^{-c}$, which lie in $U\cap U^c$.

(c) Conjugate the uniformiser and the root that $\delta$ moves; $\delta^{-1}$ moves it back, and conjugation inverts the root of unity. Inertia acts as the identity on $W$, so the cocycle is a homomorphism on inertia, factoring through the procyclic tame quotient: $\tau=\delta^m$ there, and $\tau^{-c}=(\delta^{-c})^m$.

(d) On a module killed by $\ell^N$ the cyclotomic character acts through its reduction, so equivariance is $y(gug^{-1})=\mathrm{cyc}(g)\,y(u)\circ g^{-1}$. The member is unramified at $P_0$, $\bar P_0$, so it vanishes at $\gamma_0$ and $\gamma_0^c$. If it vanishes on $H$, then for $u\in U$, $[u,\gamma_*]\in H$ gives
$0=y(u)-y(u)\circ\zeta_*^{-1}$, so $y(u)$ kills the image of $1-\zeta_*^{-1}$, which is all of $W$; and $y$ vanishes off $U$.

(e) Reciprocity at the split prime $P$ asks normalised generators at $P$ and $\bar P$; by (c), with $\delta^{-c}$ normalised at the conjugate place, pass to a multiple of $z$ for which the values at $\gamma$, $\gamma^{-c}$ are those at $\delta$, $\delta^{-c}$. The relation is then the reciprocity law that $y$ obeys at $P$, where $z$ may ramify
and $y$ is unramified.
\end{proof}

\begin{theorem}[Greenberg--Wiles count and reciprocity]\label{thm:selmer-count}
Let $J$ be tied primes with $V_i$ the triviality on $T/\ell^M$, $M<N$ and $a\mid\ell^M$.
\begin{enumerate}[label=\textup{(\alph*)}]
\item The tied Selmer group is no longer than its dual $\mathrm{Sel}^*_a(J)$, a module of relations at split primes unramified at $P_0$, $\bar P_0$, up to the places term.
\item If some member $z$ of the tied Selmer group has a value with a unit coordinate and, at every new auxiliary prime $P$ with Frobenius $\varphi$ and tame generator $\gamma$, the values of $z$ at $\varphi$, $\varphi^{-c}$ are the values at $\gamma$, $\gamma^{-c}$ of
a tied class that may ramify at $P$, $\bar P$, then the dual group has length at most the global term.
\end{enumerate}
\end{theorem}

\begin{proof}
(a) By \cref{lem:tied-setting}(d), (e) and (c), with $\#S'$ the number of primes of the setting,
\begin{align*}&\length W+\length\mathrm{Sel}_a(J)\le\length Z(\emptyset)\\
&\qquad\le\length\mathrm{Sel}^*_a(J)+\length W+\textstyle\sum_{v\in S'}\length L_v-(\#S'+f)\length W\\
&\qquad\le\length\mathrm{Sel}^*_a(J)+\length W+\text{places term},\end{align*}
and the finite length of $W$ cancels. The first prime was removed from the allowed ramification by \cref{lem:tied-setting}(a).

(b) Let $y\in\mathrm{Sel}^*_a(J)$ and put $\Psi_t(h,h')=y(h')(tz(h))-y(h^c)(tz(h'^c))$ on $H\times H$. \emph{First step: $\Psi_t(h,h)=0$.} As $z$ and $y$ vanish at $\gamma_0$, $\gamma_0^c$ and $\gamma_0$ acts as the identity, $\Psi_t$ has the same value at $(h,h)$ and at $(g,g)$, $g=\gamma_0h$. By \cref{lem:mirror}(b) take a new
auxiliary prime whose Frobenius $\varphi$ is so near $g$ that $\Psi_t(\varphi,\varphi)=\Psi_t(g,g)$, by local constancy. By hypothesis some $z'$ has $z'(\gamma)=z(\varphi)$ and $z'(\gamma^{-c})=z(\varphi^{-c})=-z(\varphi^c)$, as $\varphi^c\in U$ acts as the identity; reciprocity (\cref{lem:mirror}(e)) for $z'$ is
$y(\varphi)(tz(\varphi))-y(\varphi^c)(tz(\varphi^c))=0$, that is $\Psi_t(\varphi,\varphi)=0$. \emph{Second step: $\Psi=0$.} $\Psi_t$ is biadditive on $H$ (both $z$ and $y$ are homomorphisms on $H\subseteq U$), so it is alternating: $\Psi_t(h,h')+\Psi_t(h',h)=0$. Conjugating the first argument by $\gamma_*$ multiplies $t$ by $\zeta_*$, and conjugating the
second divides it by $\zeta_*$ (\cref{lem:selmer-basics}(c), \cref{lem:mirror}(d)); subtracting the alternation at $(h,h')$ with the parameter $\zeta_*^{-1}t$ from that at $(\gamma_*h\gamma_*^{-1},h')$ gives $\Psi_{(\zeta_*-\zeta_*^{-1})t}(h,h')=0$ for all $t$, and $\zeta_*-\zeta_*^{-1}=\zeta_*^{-1}(\zeta_*^2-1)$ is a unit. \emph{Third step.} By \cref{lem:sah}
$z(H)$ spans $W$. If $\sum_h\alpha_hz(h)=0$, then $\Psi=0$ shows that $\sum_h\alpha_hy(h^c)$ kills all $tz(h'^c)$, $h'\in H$, which span; so $\eta(z(h))=y(h^c)$ defines a linear map $\eta\colon W\to W^*$. It is equivariant for the twisted action, by the equivariance of $y$ and $gz(h)=z(ghg^{-1})$, so
$\eta\in\Hom_G(W,W^{*c})$. If $\eta=0$, then $y$ vanishes on $H^c=H$, hence vanishes (\cref{lem:mirror}(d)). Hence $y\mapsto\eta$ is injective, and the dual group has length at most the global term.
\end{proof}

\subsection{The descent}\label{ss:descent}
For a uniformiser $\varpi$ with $\varpi^e\mid\ell^M$, the \emph{total tied length} is the sum over the data and the levels $\varpi^k$, $1\le k\le e$, of the lengths of the tied Selmer groups.
\lean{New.Selmer.tilting, New.Selmer.exists_injective_quotSpan_onto_torsion, New.Selmer.isUnit_of_pow_smul_ne_zero, New.Selmer.exists_length_selmerModule_pow_le, New.Selmer.exists_unit_residue_notMem, New.Selmer.sum_lt_sum_of_ne_top, New.Selmer.exists_isTiedPrimes_singleton, New.Selmer.SelmerDatum.finite_tiedSelmer, New.Selmer.SelmerDatum.selmerModule_inf_ker_le, New.Selmer.SelmerDatum.tiedSelmer_le_of_slope_eq_zero, New.Selmer.SelmerDatum.tiedSelmer_insert_inf_ker_le, New.Selmer.SelmerDatum.tiedSelmer_update_le, New.Selmer.SelmerDatum.exists_length_tiedSelmer_insert_le, New.Selmer.SelmerDatum.length_tiedSelmer_le_of_isUnit, New.Selmer.SelmerDatum.length_tiedSelmer_pow_le, New.Selmer.totalTiedLength, New.Selmer.exists_slope_totalTiedLength_lt, New.Selmer.exists_isTiedPrimes_length_tiedSelmer_le, New.Selmer.exists_totalTiedLength_le, New.Selmer.exists_isTiedPrimes_card_le_and_length_le}

\begin{lemma}[Tools of the descent]\label{lem:descent-tools}
\begin{enumerate}[label=\textup{(\alph*)}]
\item (Tilting.) Let $\mathrm{ev}_\gamma,\mathrm{ev}_\varphi$ be linear maps from a module to $(O/I)^m$ of finite length, and $C$ a submodule of the source. For $u$ outside at most $m$ residue classes, the kernel of $u\,\mathrm{ev}_\gamma-\mathrm{ev}_\varphi$ on $C$ is no longer than the part of $\ker\mathrm{ev}_\gamma|_C$ that
$\mathrm{ev}_\varphi$ maps into $\mathrm{ev}_\gamma C$.
\item (The exponent.) For a uniformiser $\varpi$ and $k\ge1$, the Selmer module modulo $\varpi^k$ is, for some $i\le k$, no longer than the module modulo $\varpi^i$, and if $i\ne0$ the latter has a member with a value with a unit coordinate.
\item (Room for slopes.) In a local ring whose residue field has $q$ elements, fewer than $q-1$ residues miss the residue of some unit.
\item (The first prime.) For $g_0$ inertial at a place of $S_0$ prime to $\ell$ and a neighbourhood $W$ of $g_0$, there is a prime $P_0$ whose norm no prime of $S_0$ divides such that $\{P_0\}$ with all slopes $0$ is a tied set.
\item (Adding a prime.) Tied Selmer groups are finite. If $P$ is a new auxiliary prime for datum $i$, a cocycle with the ties of $J$ that may ramify also at $P$, $\bar P$ and vanishes at $\gamma$, $\gamma^{-c}$ lies in $\mathrm{Sel}_a(J)$; a member of $\mathrm{Sel}_a(J\cup\{P\})$ lies in $\mathrm{Sel}^P_a(J)$ and satisfies $\mathrm{ev}_\gamma z=s\,\mathrm{ev}_\varphi z$;
and for units $u$ outside at most $2n$ residue classes, adding $P$ with slope $u^{-1}$ gives a group no longer than the part of $\mathrm{Sel}_a(J)$ whose values at $\varphi$, $\varphi^{-c}$ are the values at $\gamma$, $\gamma^{-c}$ of a member of $\mathrm{Sel}^P_a(J)$.
\end{enumerate}
\end{lemma}

\begin{proof}
(a) Let $C'$ be the sum of all the kernels of $u\,\mathrm{ev}_\gamma-\mathrm{ev}_\varphi$ in $C$, together with the part of $\ker\mathrm{ev}_\gamma$ in question; then $\mathrm{ev}_\varphi C'\subseteq\mathrm{ev}_\gamma C'$. Write $\mathrm{ev}_\gamma C'$ as the image of a map $\pi$ on a free module $\Lambda$ of rank at most $m$ (the preimage of $\mathrm{ev}_\gamma C'$ in
$O^m$), choose $\sigma\colon\Lambda\to C'$ with $\mathrm{ev}_\gamma\circ\sigma=\pi$ and an endomorphism $U$ of $\Lambda$ with $\pi\circ U=\mathrm{ev}_\varphi\circ\sigma$. Off the at most $m$ residual roots of the characteristic polynomial of $U$, $\det(u-U)$ is a unit, and $(u\,\mathrm{ev}_\gamma-\mathrm{ev}_\varphi)\circ\sigma=\pi\circ(u-U)$
is onto $\mathrm{ev}_\gamma C'$. Hence $u\,\mathrm{ev}_\gamma-\mathrm{ev}_\varphi$ has no smaller image on $C'$ than $\mathrm{ev}_\gamma$; every map splits the length of $C'$ into image and kernel, and the image has finite length, so the kernel is no longer.

(b) Let $\varpi^i$ be the exponent of the module modulo $\varpi^k$. Multiplication by $\varpi^{k-i}$ identifies $O/\varpi^i$ with the $\varpi^i$-torsion of $O/\varpi^k$, so the module modulo $\varpi^i$ is the $\varpi^i$-torsion of the module modulo $\varpi^k$, that is, all of it. A member not killed by $\varpi^{i-1}$ has a value with a coordinate not killed by $\varpi^{i-1}$, which is a
unit of $O/\varpi^i$.

(c) The units have $q-1$ residues.

(d) Apply \cref{prop:aux-primes}(a) avoiding the primes that divide the norm of a prime of $S_0$, which include $S_0$. If $v\in S_0$ divided the prime number $\Nm P_0$, then $\Nm P_0\mid\Nm v$ and $P_0$ would divide $\Nm v$. The conditions on later primes are empty.

(e) Tied Selmer groups are modules of maps on a finite quotient of $G$ with values in the finite $W$. The second claim is \cref{prop:aux-primes}(d) at $P$: adding $P$ with slope $0$ gives a tied set whose only new prime is $P$. The third is the definition of the ties at $P$ with slope $s$. For the last, the tie with slope $u^{-1}$ says $u\,\mathrm{ev}_\gamma=\mathrm{ev}_\varphi$ on
$C=\mathrm{Sel}^P_a(J)$; by (a) with $m=2n$ (values at $\gamma$ and $\gamma^{-c}$), off at most $2n$ residues the kernel is no longer than the part of $\ker\mathrm{ev}_\gamma|_C$ that $\mathrm{ev}_\varphi$ maps into $\mathrm{ev}_\gamma C$, and $\ker\mathrm{ev}_\gamma|_C\subseteq\mathrm{Sel}_a(J)$ by the second claim.
\end{proof}

\begin{lemma}[Steps of the descent]\label{lem:descent-steps}
Let $\varpi^e\mid\ell^M$, $M<N$, let $V_i$ be the triviality on $T_i/\ell^M$, and let the residue field have more than $2n_ie+1$ elements. Call a member of $\mathrm{Sel}_a(J)$ \emph{in the tame image} at a new auxiliary prime $P$ if its values at $\varphi$, $\varphi^{-c}$ are the values at $\gamma$, $\gamma^{-c}$ of a member of
$\mathrm{Sel}^P_a(J)$. The \emph{total tied length} is the sum of the lengths of $\mathrm{Sel}_{\varpi^k}(J)$ over the data and $1\le k\le e$.
\begin{enumerate}[label=\textup{(\alph*)}]
\item If some member of $\mathrm{Sel}_a(J)$ has a value with a unit coordinate and all members are in the tame image at every new auxiliary prime, then $\mathrm{Sel}_a(J)$ has length at most the global term plus the places term; if all members modulo every $\varpi^k$, $1\le k\le e$, are in the tame image, the same holds for
$\mathrm{Sel}_{\varpi^e}(J)$.
\item (One step.) If for some $1\le k\le e$ some member of the tied Selmer group of $T_i$ modulo $\varpi^k$ is not in the tame image at a new auxiliary prime $P$, then $P$ can be added with a unit slope for $i$ and slope $0$ for the other data so that the total tied length drops.
\item There are tied primes $J'$ with the same $P_0$, at most $\#J$ plus the total tied length of $J$ in number, whose tied Selmer groups modulo $\varpi^e$ have length at most the global term plus the places term.
\item If all slopes are $0$, the total tied length has a bound independent of $N$, $P_0$, $J$ and the data at the primes.
\end{enumerate}
\end{lemma}

\begin{proof}
(a) The first claim is \cref{thm:selmer-count}(a) and (b). For the second, pass by \cref{lem:descent-tools}(b) to a level $\varpi^k$, $k\le e$, at which the group is no shorter and, unless $k=0$, has a member with a unit value, and apply the first claim there.

(b) At each level $m\le e$, \cref{lem:descent-tools}(e) excludes at most $2n_i$ residues; as the residue field has more than $2n_ie+1$ elements, \cref{lem:descent-tools}(c) gives a unit $u$ good at all levels. Adding $P$ with slope $u^{-1}$ for $i$ and $0$ for the others, no group grows: for the other data the new group lies in the old one by
\cref{lem:descent-tools}(e) (slope $0$), and for $i$ it is no longer than the part in the tame image. At the level $k$ this part is a proper submodule, and a proper submodule of a module of finite length is strictly shorter; so the sum drops.

(c) Among the tied sets $(J',x')$ with the same first prime whose number of primes plus total tied length is at most that of $(J,x)$, take one with the least total tied length. If some member at some level failed to be in the tame image at some new prime, (b) would give a set with one more prime and total length at least $1$
smaller, still in the family, contradicting minimality. Hence (a) applies.

(d) With all slopes $0$ the tied groups lie in the Selmer groups $\mathrm{Sel}(\emptyset)$ without auxiliary primes (\cref{prop:aux-primes}(d)), which are finite and depend on none of these choices.
\end{proof}

\begin{theorem}[Descent]\label{thm:descent}
Let $T_i$ be finitely many Selmer data of ranks $n_i$, unramified outside $S_0$, $g_0$ inertial at a place of $S_0$ prime to $\ell$, $W$ a neighbourhood of $g_0$, and let the residue field have more than $2n_iv(\ell^M)+1$
elements. There is an integer $s$ such that for every $N>M$ there is a set $J$ of at most $s$ tied primes whose tied Selmer groups modulo $\ell^M$ have length at most the global term plus the places term.
\end{theorem}

\begin{proof}
Let $\varpi$ be a uniformiser and $e=v(\ell^M)$, so $\varpi^e$ and $\ell^M$ generate the same ideal and the residue field has more than $2n_ie+1$ elements. Let $b$ bound the total tied length for all slopes $0$ (\cref{lem:descent-steps}(d)); it does not depend on $N$. For $N>M$ take the first prime $P_0$ of \cref{lem:descent-tools}(d), with
Frobenius element in $W$, and $J=\{P_0\}$ with slopes $0$. \Cref{lem:descent-steps}(c) gives a tied set with at most $1+b$ primes whose tied Selmer groups modulo $\varpi^e=\ell^M$ have length at most the global term plus the places term. Take $s=1+b$.
\end{proof}

\subsection{Main summands as Selmer data}\label{ss:main-summand-data}
\lean{New.Selmer.exists_dvd_ell_pow, New.Selmer.ne_top_of_normLength_le, New.Selmer.IsMainSummandLike.nonempty, New.Selmer.IsMainSummandLike.trace_lattice_conjGal, New.Selmer.IsMainSummandLike.conjGal_inv_eq_smul_one, New.Selmer.IsMainSummandLike.isUnit_of_eq_smul_one, New.Selmer.IsMainSummandLike.sum_countP_neg_le, New.Selmer.IsMainSummandLike.sum_finrank_crysCocyclesAt_le, New.Selmer.IsMainSummandLike.normLength_congruenceTerm_le, New.Selmer.IsLatticeIn.normLength_addVal_det_sub_one, New.Selmer.IsMainSummandLike.normLength_localTerm_le, New.Selmer.IsMainSummandLike.normLength_firstPrimeTerms_le, New.Selmer.IsMainSummandLike.normLength_placeTerms_add_firstPrimeTerms_le, New.Selmer.IsMainSummandLike.exists_mem_inertiaAbove_eq_smul_one_isUnit, New.Selmer.IsMainSummandLike.eventually_exists_selmerDatum}

\begin{lemma}[Terms of a main summand]\label{lem:main-summand-terms}
Let $R$ be like a main summand of rank $n$, $u$ a lattice in $R$ over $O$ whose image spans all matrices, and normalise lengths so that $O/\ell$ has length $1$.
\begin{enumerate}[label=\textup{(\alph*)}]
\item If $g$ acts through $u$ by a scalar $y$, then $g^{-c}$ acts by $\mathrm{cyc}(g)^{t-1}y$, and $y$ is a unit.
\item The negative Hodge--Tate jumps, summed over the complex embeddings, number at most $fn$, and the ranks of the crystalline cocycles at the places above $\ell$ add up to at most $n(\#\{w\mid\ell\}+f)$.
\item The global term is at most $n\,v_\ell(n_{\rm ch})$; the local term at each place above $p_0p_1$ prime to $\ell$ is at most $n\,v_\ell(n_{\rm ch})$; the two terms $v(\det(u(g)-1))$ at the tame generator $\gamma$ and $\gamma^{-c}$ add up to at most $2n\,v_\ell(n_{\rm ch})$; in all, the local terms are at most $(4f+2)n\,v_\ell(n_{\rm ch})$.
\item Some inertial element above $p_0p_1$, prime to $\ell$, acts by a scalar $\zeta$ with $\zeta^2-1$ a unit.
\end{enumerate}
\end{lemma}

\begin{proof}
(a) As the $u(x)$ span all matrices, it is enough to compare traces against them. Conjugate self-duality on traces, $\tr u(g^c)=\mathrm{cyc}(g)^{-(t-1)}\tr u(g^{-1})$, applied at $x^cg^{-1}$ and at $x^c$, makes both traces multiples of $\tr u((x^c)^{-1})$ with the claimed ratio. A scalar of a lattice element is a unit, as $u(g)$ is invertible over $O$.

(b) An embedding and its conjugate have at most $n$ negative jumps together (\cref{def:main-summand}), and there are $f$ pairs. The rank at $v\mid\ell$ is $n$ plus the number of negative jumps at the embeddings of $v$ (\cref{thm:B-padic}(a)), and each embedding is counted at one place at most.

(c) Take $g$ acting by a scalar $\alpha$ with cyclotomic character $p_l^n$ (\cref{def:main-summand}). By (a), $\alpha-\mathrm{cyc}(g^c)\alpha'=(1-p_l^d)\alpha$ with $d=nt$ ($d=n$ if $t=0$), $\alpha$ a unit and $1\le d\le\mathrm r\mathrm t$; so $v(\alpha-\mathrm{cyc}(g^c)\alpha')=v_\ell(p_l^d-1)\le v_\ell(n_{\rm ch})$ (\cref{lem:lift-settings}(b)), and
\cref{lem:sah} gives the global term. At a place $v$ above $p_0p_1$ prime to $\ell$, a power $m$ prime to $v$ kills the image of inertia and $u$ is trivial on inertia near $1$ (\cref{def:main-summand}). A local cocycle is determined by its value at the tame generator together with its value at a Frobenius $\varphi$ when that first value vanishes: such a cocycle is
unramified, and $\gamma\varphi=\varphi(\varphi^{-1}\gamma\varphi)$ with both $\gamma$ and $\varphi^{-1}\gamma\varphi$ inertial shows that $u(\gamma)$ fixes its value at $\varphi$. Hence the local cocycles have length at most $\length W+\length\ker(u(\gamma)-1)$ for the inertial $\gamma$ of the row, and the local term is at most
$v(\det(u(\gamma)-1))$, which is at most $n\,v_\ell(n_{\rm ch})$ by \cref{def:main-summand}; the valuation of $\det(u(g)-1)$ in $O$, normalised, is that of $\det(R(g)-1)$ in $\bar\Q_\ell$. At most $4f$ places above $p_0p_1$ are prime to $\ell$, and the tame generator gives two more terms.

(d) The scalar $y$ of \cref{def:main-summand}, with $y^{\ell^aj}-1$ a unit for $1\le j\le4r+2$, lies in $O$, being a diagonal entry of $u(\gamma)$, a conjugate of $R(\gamma)$; take $j=2$, $a=0$.
\end{proof}

\begin{proposition}[Main summands are Selmer data]\label{prop:main-summand-data}
Let $R$ be like a main summand and $\gamma$ the tame generator at a chosen place prime to $\ell$. Over all large coefficient rings there is a Selmer datum $T$ for $\gamma$ and the bad places, a lattice in $R$, with places those above
$\ell p_0p_1$, local relations those of the crystalline cocycles above $\ell$, and $\beta_T$ of normalised length at most $(4f+3)\,n\,v_\ell(n_{\rm ch})$.
\end{proposition}

\begin{proof}
Take the lattice $u$ of \cref{lem:main-summand-terms} and the element $\gamma_*$ of \cref{lem:main-summand-terms}(d); the cyclotomic character is trivial at the inertial $\gamma_*$, so by \cref{lem:main-summand-terms}(a) $\gamma_*^{-c}$ acts by the same scalar $\zeta_*$. The places are those above $\ell p_0p_1$, outside which $R$ is unramified; the
modules $Z_w$ are the crystalline cocycles above $\ell$, finitely generated with the rank bound of \cref{lem:main-summand-terms}(b); Burnside's form of irreducibility is the spanning of the image. By \cref{lem:main-summand-terms}(c), the global term is at most $n\,v_\ell(n_{\rm ch})$ and the local terms at most
$(4f+2)n\,v_\ell(n_{\rm ch})$; so $\beta_T$ has normalised length at most $(4f+3)n\,v_\ell(n_{\rm ch})$, and in particular all terms are finite.
\end{proof}
